% !TeX root = ../papers/nc-p3.tex
\section{負交換の射影商、周期一の因子と退化した楕円枠}
\label{sec:negative-pairing-unmarked-quotient}
定理\ref{thm:negative-component-elliptic-graph}の負交換の座標を用い、
$zD(z^2-1)\ne0$, $D=ad-bc$ とする。$A=[a:b:c:d]\in\mathbb P^3$ であり、
正則点とは定理\ref{thm:exceptional-exchange-eight-lines}の八直線を
$(a,-b,-c,d)$ に対して除いた点をいう。以下
\[
 p=a+d,\quad q=a-d,\quad r=zb+c,\quad t=zb-c
\]
とする。
\begin{theorem}\label{thm:negative-pairing-unmarked-quotient}
$qt\ne0$ の全正則点で、指定始点の同型類は次の次数16の群 $H$ の軌道である。
$U,V,S$ はそれぞれ $z$ を $1/z$ とし、射影四座標を
\[
 U:[a:c:b:d],\qquad V:[d:zb:c/z:a],\qquad S:[a:zb:c/z:d]
\]
に移す。第四生成元 $B$ と、一位置移動 $J$ は $z$ を保ち、$(p,r,q,t)$ には
\[
 B:\big(2zp+(z+1)r,-z(z+1)p-2zr,-(z-1)t,-z(z-1)q\big),
 \qquad J:(r,zp,-t,zq)
\]
で作用する。$z$ を固定する部分群は位数8の二面体群である。
この全域で最小周期は $r=-p$ または $r=-zp$ のときちょうど1、その他で2である。
前二平面は商上で一つの既約因子を与える。一般商は有理四次元である。
\end{theorem}

\begin{proof}[全正則点での頂点の復元]
$X_i=x_i^2$, $P=D(X_0-X_1)-(a+b+c+d)X_2+(az+bz^2+c+dz)X_3$,
$Q=aX_0-dX_1-X_2+zX_3$ とする。左の四次完全交叉 $E$ は $P=Q=0$ である。
$Dq\ne0$ により前節の二本の二次式と同じイデアルを張る。
平方座標の直線上で $(X_0,X_1)$ と $(X_2,X_3)$ の各対は独立である。
実際、前節の $LO-MN=(a-d)(bz-c)(z-1)D$ である。
各座標に従う二重被覆の分岐点となる四つの零点は、各対の中では重ならない。
全て異なれば $E$ は滑らかである。重複が一対なら、二文字の平方の比例を因数分解して
二つの滑らかな平面二次曲線を得る。これらは他の二文字が定める座標直線上の二点で横断する。
二対なら二つの独立な平方比例を因数分解し、被約な四直線の四角形を得る。
重複はいずれも $\{0,1\}$ と $\{2,3\}$ から一文字ずつの対である。

全ての四次小行列式は $E$ 上で消える。これは関数体上での割算を
$X_0,X_1$ について行う際の分母が $Dq$ の冪だけであることから特殊化できる。
$E$ 外では点は前節の四孤立点に限ることを示す。行を0から数え、$m_I$ を四次小行列式とする。
$m_{0145},m_{2345}$ の和と差から $x_2x_3 Q=x_0x_1 P=0$。
$Q=0,P\ne0$ なら $m_{0124},m_{0125},m_{0134},m_{0135}$ は
$x_1x_2,x_1x_3,zx_0x_3,x_0x_2$ に $P$ を掛けたものである。
よって片方の座標対は零である。$x_2=x_3=0$ の場合は $x_0x_1=0$ により座標頂点になるが、
$D\ne0$ で両頂点の評価は階数4である。もう一方なら $X_2=zX_3$ となる。

$Q\ne0$ の場合、$x_3=0,x_2\ne0$ なら小行列式から
$(a+b)x_0=(a+c)x_0=(c+d)x_1=(b+d)x_1=0$ となる。
一組がともに零なら正表示の八直線の左側02または12であり、除外済みである。
よって $x_0=x_1=0$ となり $m_{1245}=-Q(b(X_0-X_1)+X_2-X_3)$ に反する。
$x_2=0,x_3\ne0$ でも対応する四係数は $a+bz,az+c,c+dz,bz+d$ であり、
八直線の左側03、13を除くと同じ矛盾を得る。
残りの $x_2=x_3=0$ では $x_0x_1P=0$ と座標頂点の排除から $X_0=X_1$ である。
得た二対の点はいずれも実際の零点で、一方では $Q\ne0$、他方では $P\ne0$ であり、$E$ と交わらない。
\path{computations/verify_negative_smooth_pairing.py} に小行列式と両側の対角鉛筆のPlücker座標を与えた。

滑らかな場合の枠は前節の四頂点である。二次曲線二本なら鉛筆には階数2の曲面一つと
階数3の曲面二つがある。前者の頂点直線は重なった固有値の二頂点を結び、後二者の頂点は
残る固有値の各一頂点である。四つの孤立点の六割線のうち、後二頂点のどちらかを含むのは
座標直線 $01,23$ だけである。この二直線と頂点直線の交点が残りの二頂点である。
四角形なら相対する節点の二割線が重複固有値の二直線である。
孤立点の六割線のうちその両方と交わるのは $01,23$ だけである。実際、異なる組の
孤立点を結ぶ割線では、各組に残る一座標ずつを同時に零にできない。
四つの交点から元の四頂点を取り戻す。右と次辺については以下の $J$ の恒等式が
$qt\ne0$ と正則性を保つので同じである。この恒等式の検証は頂点復元を使わない。
\end{proof}

\begin{proof}[全置換と薄い係数の条件]
非零台には、上位ビットが0の文字と1の文字が二つずつある48語が常に残る。
四文字の各置換は $\mathbb F_2^2$ 上アフィンである。総和零の64語のうち48語で
変換後の総和が零なら、全64語で零である（非恒等なアフィン線型条件の解は高々32語）。
従って置換は四位置で共通の線型部 $\phi$ と総和零の移動量を持つ。
各隣接切断の厚い色が $0,1$ なので、$\phi$ は恒等または $2,3$ の交換で、
四移動量の上位ビットは等しくなければならない。各位置移動につき32候補である。

48語の係数は $\pm z^e$ であり、$z=w_{0022}w_{1133}/(w_{0033}w_{1122})$。
候補の変換後のこの比 $z'$ は $z$ または $1/z$ である。元の表示を $z'$ とした薄い係数と
候補側の薄い係数の比について、文字倍率行列の飽和整数核で符号・指数を比較する。
移動が $k$、各文字移動が $u_i$ のとき $u_0+u_2=k\bmod2$ の16候補では全て恒等式、
他の16候補の差のLaurent多項式の最大公約数は $z-1$ である。
薄い係数を固定した後、左上 $A$ の四成分の倍率は等しい。対応する四語の任意の二つの
文字指数の差は薄い48語の整数行格子に入る。従ってこの主張は $A$ に零成分があっても成り立つ。
共通倍率は全語に対する一体の変更で実現できる。各語で上位ビットの和を $n$ とすると
共通倍率 $v$ に対する因子は $v^{1-n/2}$ である。

指定始点で残る置換は、移動 $(0,1,0,1),(1,0,1,0)$、線型部の交換、および
この交換に移動 $(2,2,2,2)$ を加えたものにより生成される。これが順に $U,V,S,B$ である。
一位置移動にはさらに $(0,0,1,1)$ の文字移動を選ぶと $J$ となる。
これら五変換に対して実際の全64係数の比は、符号と $z,z-1,D$ の整数冪であり、
その全ての文字格子関係が成立する。よって開条件上の各点で対角倍率が存在し、
係数が消える場合にも単元を掛けた恒等式を特殊化できる。
\path{computations/verify_negative_component_quotient.py} が五つの作用と薄い場合分けを検証する。
\end{proof}

\begin{proof}[周期、安定群、有理商]
$US,VS$ は $(q,t)$ の二つの符号を独立に変え $(p,r)$ を保つ。
$B$ の最初の行列を $B_0$、$J$ の最初の行列を $J_0$ とする。
$B_0^2=-z(z-1)^2 I$, $J_0^2=zI$, $B_0J_0=-J_0B_0$ である。
従って $B^2$ は射影的に $(q,t)$ の同時符号変更で、$B$ は二つの単一符号変更を入れ替える。
これで $z$ を固定する八元の二面体群 $H_0$ を得る。$U$ は $z$ を反転し、$H$ の全16元を尽くす。
$J^2$ は二つの符号交換の積を含む $H_0$ の元である
（直接には $(p,r,q,t)\mapsto(p,r,-q,-t)$）。従って常に周期2以下である。

周期1なら $Jx=h x$, $h\in H_0$ が射影的に成立する。
$h$ が符号変更だけの場合、$(p,r)\ne(0,0)$ なら $r^2=zp^2$、さらに $qt\ne0$ から
$t^2=zq^2$ となる。これは $4D=p^2-q^2-(r^2-t^2)/z\ne0$ に反する。
$(p,r)=(0,0)$ は二平面の共通部分である。
残る四元では $J_0(p,r)=cB_0(p,r)$ と書き、後の二座標を比較する。
$qt\ne0$ により $c=\pm1/(z-1)$ で、二つの符号は互いに逆である。
各符号で最初の二式は $r=-zp$ または $r=-p$ と同値で、逆に各平面で必要な二符号を選べる。
二平面は $B$ で入れ替わる。各平面には正則点が稠密である
（除外した八直線と $qtD=0$ はいずれも平面全体を含まない）ので商上の一因子である。

安定群も全点で数えられる。台の文字指数行列は薄い48行と左上一語だけで
Smith因子 $1^{11},2,2$ を持ち、全64語の行格子に等しい。
$D\ne0$ によりそのような左上一語があり、どの左上一語を選んでも同じである。
従って射影安定群の位数は $4|\operatorname{Stab}_{H_0}(x)|$ である。$(p,r)\ne(0,0)$
では後の因子は通常1で、$t^2=-zq^2$ かつ $(p,r)$ が $B_0$ の固有線にあるときだけ2になる。
実際、単一の符号変更では固定できず、$B$ を含む元の二符号はその固有値によって決まる。
$(p,r)=(0,0)$ では同時符号変更も効かず、後の因子は2、ただし $t^2=-zq^2$ なら4である。
$t^2=zq^2$ はここでは $D=0$ で除かれている。群そのものは該当する置換の文字倍率の解を合成すればよい。

最後に関数体を直接求める。二つの独立な符号と $U$ の不変体を
\[
 Z=(z+1)^2/z,\quad x=r/((z+1)p),\quad u=q^2/p^2,\quad v=t^2/(zp^2)
\]
が生成する。実際、先に符号を除いた後は $z$ を戻すだけで元の体に戻る。
$B$ は $x'=-(1+2x)/(2+Zx)$, $\delta=(Z-4)/(2+Zx)^2$ として
$(x,u,v)\mapsto(x',\delta v,\delta u)$ とする。従って
\[
  Z,\quad x+x',\quad (1+\delta)(u+v),\quad
  (1+\delta)(u-v)/(x-x')
\]
は不変である。$x+x'$ は $x$ の真の二次有理式なので、この四元に $x$ を戻せば
$u,v$ まで解け、拡大次数は2である。これで四つの独立な生成元を得た。
\end{proof}
